\scanpage{147}
% lines 1–4
Soit $C$ une courbe complète, réduite et connexe sur $k$ \inserted{alg. clos}

On veut calculer $\mathrm{Pic}_C$
\[
\unsure{\cdot|0}\unsure{-|\to}\underline O_C^{*}-\underline{\widetilde O}_C^{*}\to\prod_P A_P^{\unsure{*}}-0\qquad P\ \text{les pts doubles}
\]
\[
\cdot\to\underbrace{G_m-\deleted{\unsure{P}}\prod_{C_i}G_m\to\prod_{P_\alpha}A_P}_{\text{\unsure{gr.|tor.|fin|tore}}}\to\mathrm{Pic}_{C/k}\to\deleted{\unsure{\mathrm{Pic}}}\prod\mathrm{Pic}_{\widetilde C_i/k}\to0
\]
\margin{$k[[t]]/t^2$}
\figure{\begin{tikzpicture}[x=.65cm,y=.65cm]
\fill (0,0) circle (.075); \draw[thick] (0,0) .. controls (.55,0) and (.88,.11) .. (1.15,.45); \draw[thick] (0,0) .. controls (.65,-.01) and (.83,-.15) .. (1.08,-.39);
\end{tikzpicture}}

% lines 5–9
On trouve que l'on a \deleted{dans} $\mathrm{Pic}_{C^{0}/k}$ une partie linéaire, \deleted{et \unsure{en fait} multiplicative,} qui est isomorphe \deleted{à}
\[
G=\left(\prod_{P_\alpha}A_{P_\alpha}\right)\big/\left(\prod G_m/G_m\right)
\]
donc est de dimension égale à
\[
\dim G=\sum\mu_\alpha-c_0+1
\]
% slanted
\margin{$c_0=$ nb des comp. de $C_i$}
% lines 10–11
où $\mu_\alpha=\dim P_\alpha\inserted{=\dim_k\underline O_{\widetilde C}/\underline O_C}$ \deleted{\unsure{pour simplifier} \unsure{en} $P_\alpha$}. Posons $\mu_\alpha=1+\delta_\alpha$, \deleted{on trouve donc}

% lower cancelled block, lines 12–18
\deleted{\[
\dim G=c_1-c_0+1+\sum\delta_\alpha=b_1+\sum\delta_\alpha
\]
\deleted{\unsure{de}} où $b_1$ est le nb de cycles du schéma combinatoire de $C$.}
\margin{$c_0-c_1=1-b_1$}
\deleted{La dimension de la partie \deleted{multiplicative} \inserted{additive} de $G$ est
\[
\begin{aligned}
&\deleted{\sum_\alpha\unsure{(\nu_\alpha-1)}=\left(\sum\unsure{\nu_\alpha}\right)-c_1=2c_0-c_1=}\\
&\deleted{\sum_\alpha(\mu'_\alpha-\unsure{1})-c_0+1}\qquad\sum(\mu_\alpha-\mu'_\alpha)
\end{aligned}
\]
où $\mu'_\alpha$ est la multiplicité ``ensembliste'' de $P_\alpha$, $=$ nb de \deleted{\unsure{courbes}} \inserted{branches \unsure{formelles|simples}} passant \deleted{\unsure{e}} par $\alpha$.}
% archive sheet number: 146

